% SPDX-FileCopyrightText: © 2026 Idiap Research Institute <contact@idiap.ch>
%
% SPDX-FileContributor: Andre Freitas <andre.freitas@idiap.ch>
% Neuro-symbolic AI Group
%
% SPDX-FileContributor: Danilo Gusicuma <danilo.gusicuma@idiap.ch>
% Neuro-symbolic AI Group
%
% SPDX-License-Identifier: MIT
% Generated from step-explanations.json
\paragraph{Step-by-step reconstruction.}
A proposed Yang--Mills programme claims that decay of finitely many Wilson-loop observables establishes the full mass gap.
The following is a source-guided illustrative reconstruction. Internal procedures are prescriptions; the native declarations do not establish their execution.
\begin{enumerate}[leftmargin=*,label=\arabic*.]
\item \textbf{Identify the programme} (\texttt{scope\_\allowbreak{}activation}).
\begin{description}[style=standard,labelindent=0pt,leftmargin=5.5em,labelwidth=5em,labelsep=0.5em,align=left,font=\normalfont\bfseries,itemsep=0.15em,parsep=0pt,topsep=0.3em]
\item[Purpose] Separate the scientific target from its support.
\item[Input] Activated domain target; intended claim
\item[Method] The native scope producer is unresolved. Target-freezing instructions distinguish existence, spectral exclusion, clustering, confinement and area law before selecting a programme.
\item[Check] Different endpoints require different premises and cannot be merged by a shared label.
\item[Output] Intended activation: pure-gauge programme; separate target burdens
\item[Next] Intake calibration consumes the activation; its existence does not establish any target theorem.
\item[Example] Split the announcement into finite-observable decay, observable completeness, continuum existence, reconstruction and spectral exclusion instead of accepting one as all the others.
\end{description}
\item \textbf{Judge whether inquiry can begin} (\texttt{intake\_\allowbreak{}disposition}).
\begin{description}[style=standard,labelindent=0pt,leftmargin=5.5em,labelwidth=5em,labelsep=0.5em,align=left,font=\normalfont\bfseries,itemsep=0.15em,parsep=0pt,topsep=0.3em]
\item[Purpose] Assess whether the scoped question is usable.
\item[Input] Scope activation; domain; target
\item[Method] Assess whether the target and its boundaries are definite enough to investigate, retaining missing claims as obligations.
\item[Check] Readiness concerns a usable question, not support for its ultimate conclusion.
\item[Output] Intake disposition: ready, blocked or stalled
\item[Next] Original policy routes ready intake to grounding; the native return constructs no grounding request.
\item[Example] Proceed only with the separate obligations visible; an unspecified meaning of mass gap requires target clarification before scientific grounding.
\end{description}
\end{enumerate}
\paragraph{Changed-input contrast.}
If intake is instead asked only for a finite-regulator observable estimate, continuum reconstruction is not part of that narrower result. The activation must preserve the selected target rather than silently exchanging the two questions.
